\chapter{Long Time Behaviors}
\label{chap:cz-long-time-behavior}
\section{The Ricci Flow on Two-manifolds}
\begin{lemma} [normalized evolution of homogeneous quantities]
  \label{lem:cz-normalized-evolution-homogeneous-quantities}
  \source{cao-zhu}
  \dcref{Lemma 5.1.1}
  \group{Normalized Ricci Flow on Surfaces}

Suppose $P$ satisfies
$$
\frac{\partial P}{\partial t}=\Delta P+Q
$$
for the unnormalized Ricci flow, and $P$ has degree $k$. Then $Q$
has degree $k-1$, and for the normalized Ricci flow,
$$
\frac{\partial \tilde{P}}{\partial \tilde{t}}=\tilde{\Delta}
\tilde{P}+\tilde{Q}+\frac{2}{n}k\tilde{r}\tilde{P}.
$$
\end{lemma}
\begin{lemma} [evolution of the curvature-potential energy density]
  \label{lem:cz-curvature-potential-density-evolution}
  \source{cao-zhu}
  \dcref{Lemma 5.1.2}
  \group{Normalized Ricci Flow on Surfaces}

The function $h$ satisfies the evolution equation \begin{equation}
\frac{\partial h}{\partial t}=\Delta h-2|M_{ij}|^2+rh.
%\eqno(5.1.5)
\end{equation}
\end{lemma}
\begin{proposition} [immortality of normalized surface Ricci flow]
  \label{prop:cz-normalized-surface-flow-immortal}
  \source{cao-zhu}
  \dcref{Proposition 5.1.3}
  \group{Normalized Ricci Flow on Surfaces}

For any initial metric on a compact surface, the normalized Ricci
flow $(5.1.1)$ has a solution for all time.
%$$\eqno \#$$
\end{proposition}
\begin{theorem}[convergence on surfaces of negative Euler characteristic]
  \label{thm:cz-negative-euler-surface-convergence}
  \source{cao-zhu}
  \dcref{Theorem 5.1.4}
  \group{Normalized Ricci Flow on Surfaces}

On a compact surface with $\chi(M)<0$, for any initial metric the
solution of the normalized Ricci flow $(5.1.1)$ exists for all
time and converges in the $C^{\infty}$ topology to a metric with
negative constant curvature.
\end{theorem}
\begin{theorem}[convergence on surfaces of zero Euler characteristic]
  \label{thm:cz-zero-euler-surface-convergence}
  \source{cao-zhu}
  \dcref{Theorem 5.1.5}
  \group{Normalized Ricci Flow on Surfaces}

On a compact surface with $\chi(M)=0$, for any initial metric the
solution of the normalized Ricci flow $(5.1.1)$ exists for all
time and converges in $C^\infty$ topology to a flat metric.
\end{theorem}
\begin{lemma}[space-time Harnack estimate on positive-Euler surfaces]
  \label{lem:cz-positive-euler-surface-harnack}
  \source{cao-zhu}
  \dcref{Lemma 5.1.6}
  \group{Normalized Ricci Flow on Surfaces}

There exists a positive constant C depending only on the initial
metric such that for any $x_1$, $x_2\in M$ and $t_2>t_1\ge 0$,
$$
R(x_1,t_1)-s(t_1) \le
e^{\frac{\Delta}{4}+C(t_2-t_1)}(R(x_2,t_2)-s(t_2))
$$
where
$$
\Delta=\inf\left\{\int_{t_1}^{t_2}|\dot{\gamma}(t)|^2_tdt\ |\
\gamma:[t_1,t_2]\rightarrow M\;with
\;\gamma(t_1)=x_1,\gamma(t_2)=x_2\right\}.
$$
%$$\eqno \#$$ \vskip 0.3cm
\end{lemma}
\begin{proposition} [uniform scalar-curvature bounds on positive-Euler surfaces]
  \label{prop:cz-positive-euler-surface-curvature-bounds}
  \source{cao-zhu}
  \uses{lem:cz-positive-euler-surface-harnack}
  \dcref{Proposition 5.1.7}
  \group{Normalized Ricci Flow on Surfaces}

Let $(M,g_{ij}(t))$ be a solution of the normalized Ricci flow on
a compact surface with $\chi(M)>0$. Then there exist a time
$t_0>0$ and a positive constant $C$ such that the estimate
$$
C^{-1}\le R(x,t)\le C
$$
holds for all $x\in M$ and $t\in [t_0,+\infty).$
\end{proposition}
\begin{lemma} [decay evolution for the trace-free Hessian]
  \label{lem:cz-tracefree-hessian-evolution}
  \source{cao-zhu}
  \dcref{Lemma 5.1.8}
  \group{Normalized Ricci Flow on Surfaces}

We have \begin{equation} \frac{\partial}{\partial t}|M_{ij}|^2
=\Delta|M_{ij}|^2-2|\nabla_k M_{ij}|^2-2R|M_{ij}|^2,
\; \text{ on }\; M\times[0,+\infty). %\eqno(5.1.23)
\end{equation}
\end{lemma}
\begin{proposition}  [exponential decay of the trace-free Hessian]
  \label{prop:cz-tracefree-hessian-exponential-decay}
  \source{cao-zhu}
  \dcref{Proposition 5.1.9}
  \group{Normalized Ricci Flow on Surfaces}

Let $(M,g_{ij}(t))$ be a solution of the normalized Ricci flow on
a compact surface with $\chi(M)>0$. Then there exist positive
constants $c$ and $C$ depending only on the initial metric such
that
$$
|M_{ij}|^2\leq Ce^{-ct},\quad \text{ on }\; M\times[0,+\infty).
$$
%$$\eqno \#$$
\end{proposition}
\begin{proposition} [rigidity of compact shrinking surface solitons]
  \label{prop:cz-compact-shrinking-surface-soliton-rigidity}
  \source{cao-zhu}
  \dcref{Proposition 5.1.10}
  \group{Normalized Ricci Flow on Surfaces}

On a compact surface there are no shrinking Ricci solitons other
than constant curvature.
\end{proposition}
\begin{theorem}[convergence on surfaces of positive Euler characteristic]
  \label{thm:cz-positive-euler-surface-convergence}
  \source{cao-zhu}
  \dcref{Theorem 5.1.11}
  \group{Normalized Ricci Flow on Surfaces}

On a compact surface with $\chi(M)>0$, for any initial metric, the
solution of the normalized Ricci flow $(5.1.1)$ exists for all
time, and converges in the $C^\infty$ topology to a metric with
positive constant curvature.
%$$\eqno \#$$\vskip 1cm
\end{theorem}
\section{Differentiable Sphere Theorems in 3-D and 4-D}
\begin{theorem}[three-dimensional positive-Ricci sphere theorem]
  \label{thm:cz-three-manifold-positive-ricci-sphere}
  \source{cao-zhu}
  \dcref{Theorem 5.2.1}
  \group{Differentiable Sphere Theorems}

A compact three-manifold with positive Ricci curvature must be
diffeomorphic to the three-sphere $\mathbb{S}^3$ or a quotient of
it by a finite group of fixed point free isometries in the
standard metric.
\end{theorem}
\begin{theorem}[four-dimensional positive-curvature-operator sphere theorem]
  \label{thm:cz-four-manifold-positive-curvature-operator-sphere}
  \source{cao-zhu}
  \dcref{Theorem 5.2.2}
  \group{Differentiable Sphere Theorems}

A compact four-manifold with positive curvature operator is
diffeomorphic to the four-sphere $\mathbb{S}^4$ or the real
projective space $\mathbb{RP}^4$.
\end{theorem}
\begin{theorem}[classification under nonnegative Ricci or curvature operator]
  \label{thm:cz-nonnegative-curvature-low-dimensional-classification}
  \source{cao-zhu}
  \dcref{Theorem 5.2.3}
  \group{Differentiable Sphere Theorems}
 \
\begin{itemize}
\item[(i)] A compact three-manifold with nonnegative Ricci
curvature is diffeomorphic to $\mathbb{S}^3$, or a quotient of one
of the spaces $\mathbb{S}^3$ or $\mathbb{S}^2\times\mathbb{R}^1$
or $\mathbb{R}^3$ by a group of fixed point free isometries in the
standard metrics. \item[(ii)] A compact four-manifold with
nonnegative curvature operator is diffeomorphic to $\mathbb{S}^4$
or $\mathbb{CP}^2$ or $\mathbb{S}^2\times \mathbb{S}^2$, or a
quotient of one of the spaces $\mathbb{S}^4$ or $\mathbb{CP}^2$ or
$\mathbb{S}^3\times\mathbb{R}^1$ or $\mathbb{S}^2\times
\mathbb{S}^2$ or $\mathbb{S}^2\times \mathbb{R}^2$ or
$\mathbb{R}^4$ by a group of fixed point free isometries in the
standard metrics.
\end{itemize}
\end{theorem}
\begin{lemma} [preservation of three-dimensional Ricci pinching]
  \label{lem:cz-three-dimensional-ricci-pinching-preserved}
  \source{cao-zhu}
  \dcref{Lemma 5.2.4}
  \group{Differentiable Sphere Theorems}

For any $\varepsilon\in[0,\frac{1}{3}]$, the pinching condition
$$
R_{ij}\geq0\ \ \ and\ \ \ R_{ij}\geq\varepsilon Rg_{ij}
$$
is preserved by the Ricci flow.
\end{lemma}
\begin{proposition} [asymptotic roundness under positive Ricci curvature]
  \label{prop:cz-three-dimensional-positive-ricci-asymptotic-roundness}
  \source{cao-zhu}
  \uses{lem:cz-three-dimensional-pinching-curvature-gradient, lem:cz-three-dimensional-ricci-pinching-preserved}
  \dcref{Proposition 5.2.5}
  \group{Differentiable Sphere Theorems}

Suppose that the initial metric of the solution to the Ricci flow
on $M^3\times [0,T)$ has positive Ricci curvature. Then for any
$\varepsilon>0$ we can find $C_\varepsilon<+\infty$ such that
$$
\left|R_{ij}-\frac{1}{3}Rg_{ij}\right|\leq\varepsilon
R+C_\varepsilon
$$
for all subsequent $t\in [0, T)$.
\end{proposition}
\begin{proposition}[invariant four-dimensional curvature pinching set]
  \label{prop:cz-four-dimensional-invariant-curvature-pinching-set}
  \source{cao-zhu}
  \dcref{Proposition 5.2.6}
  \group{Differentiable Sphere Theorems}

If we choose successively positive constants $G$ large enough, $H$
large enough, $\delta$ small enough, $J$ large enough,
$\varepsilon$ small enough, $K$ large enough, $\theta$ small
enough, and $L$ large enough, with each depending on those chosen
before, then the closed convex subset $X$ of
$\{M_{\alpha\beta}\geq 0 \}$ defined by the inequalities
\begin{itemize}
\item[(1)] $(b_2+b_3)^2\leq Ga_1c_1,$ \item[(2)] $a_3\leq Ha_1 \ \
\text{and}\ \ c_3\leq Hc_1,$ \item[(3)] $(b_2+b_3)^{2+\delta}\leq
Ja_1c_1(a-2b+c)^{\delta},$ \item[(4)]
$(b_2+b_3)^{2+\varepsilon}\leq Ka_1c_1,$ \item[(5)] $a_3\leq
a_1+La_1^{1-\theta}\ \ \text{and} \ \ c_3\leq
c_1+Lc_1^{1-\theta},$
\end{itemize}
is preserved by {\rm ODE} $(5.2.3)$. Moreover every compact subset
of $\{M_{\alpha\beta}>0\}$ lies in some such set $X$.
\end{proposition}
\begin{corollary} [four-dimensional curvature-operator asymptotic roundness]
  \label{cor:cz-four-dimensional-curvature-operator-roundness}
  \source{cao-zhu}
  \dcref{Corollary 5.2.7}
  \group{Differentiable Sphere Theorems}

Suppose that the initial metric of the solution to the Ricci flow
on a compact four-manifold has positive curvature operator. Then
for any $\varepsilon>0$ we can find positive constant
$C_\varepsilon<+\infty$ such that
$$
|{\mathop{Rm}^ \circ}|\le\varepsilon R+C_\varepsilon
$$
for all $t\ge 0$ as long as the solution exists, where ${\mathop
{Rm}\limits^ \circ}$ is the traceless part of the curvature
operator.
%$$\eqno \#$$ \vskip 0.3cm
\end{corollary}
\begin{remark}[subsequential convergence in the differentiable sphere theorems]
  \label{rem:cz-differentiable-sphere-subsequential-convergence}
  \source{cao-zhu}
  \uses{thm:cz-three-manifold-positive-ricci-sphere, thm:cz-four-manifold-positive-curvature-operator-sphere}
  \dcref{Remark 5.2.8}
  \group{Differentiable Sphere Theorems}

The proofs of Theorem~\ref{thm:cz-three-manifold-positive-ricci-sphere} and Theorem~\ref{thm:cz-four-manifold-positive-curvature-operator-sphere} also show that the
Ricci flow on a compact three-manifold with positive Ricci
curvature or a compact four-manifold with positive curvature
operator is subsequentially converging (up to scalings) in the
$C^{\infty}$ topology to the same underlying compact manifold with
a metric of positive constant curvature. Of course, this
subsequential convergence is in the sense of Hamilton's
compactness theorem (Theorem~\ref{thm:cz-hamilton-compactness-for-ricci-flows}) which is also up to the
pullbacks of diffeomorphisms. Actually in \cite{Ha82} and
\cite{Ha86}, Hamilton obtained the convergence in the stronger
sense that the (rescaled) metrics converge (in the $C^{\infty}$
topology) to a constant (positive) curvature metric.
%\vskip 0.3cm
\end{remark}
\begin{theorem}[higher-dimensional pointwise-pinched sphere theorem]
  \label{thm:cz-higher-dimensional-pointwise-pinched-sphere}
  \source{cao-zhu}
  \dcref{Theorem 5.2.9}
  \group{Differentiable Sphere Theorems}

Let $n\ge4$. Suppose $M$ is a complete n-dimensional manifold with
positive and bounded scalar curvature and satisfies the
pointwisely pinching condition
$$
|W|^2+|V|^2\le\delta_n(1-\varepsilon)^2|U|^2,
$$
where $\varepsilon>0,\delta_4=\frac{1}{5},\delta_5=\frac{1}{10}$,
and
$$
\delta_n=\frac{2}{(n-2)(n+1)},n\ge 6.
$$
Then $M$ is diffeomorphic to the sphere $\mathbb{S}^n$ or a
quotient of it by a finite group of fixed point free isometries in
the standard metric.
%$$\eqno \#$$ \vskip 0.3cm
\end{theorem}
\begin{theorem}[flatness from Ricci pinching in the noncompact three-dimensional case]
  \label{thm:cz-noncompact-three-manifold-ricci-pinching-flatness}
  \source{cao-zhu}
  \dcref{Theorem 5.2.10}
  \group{Differentiable Sphere Theorems}

Let $M$ be a three-dimensional complete noncompact Riemannian
manifold with bounded and nonnegative sectional curvature. Suppose
$M$ satisfies the following Ricci pinching condition
$$
R_{ij}\ge \varepsilon Rg_{ij},\quad \text{on }\; M,
$$
for some $\varepsilon>0$. Then $M$ is flat.
%$$\eqno \#$$ \vskip 0.3cm
\end{theorem}
\section{Nonsingular Solutions on Three-manifolds}
\begin{definition}[nonsingular normalized Ricci-flow solution]
  \label{def:cz-nonsingular-normalized-ricci-flow}
  \source{cao-zhu}
  \dcref{Definition 5.3.1}
  \group{Nonsingular Three-Dimensional Solutions}

A {\bf nonsingular solution}\index{solution!nonsingular} of the
Ricci flow is one where the solution of the normalized flow exists
for all time $0\leq t<\infty$, and the curvature remains bounded
$|Rm|\leq C<+\infty$ for all time with some constant $C$
independent of $t$.
\end{definition}
\begin{theorem}[Hamilton--Ivey pinching estimate]
  \label{thm:cz-hamilton-ivey-pinching-estimate}
  \source{cao-zhu}
  \dcref{Theorem 5.3.2}
  \group{Nonsingular Three-Dimensional Solutions}

Suppose we have a solution to the $($unnormalized$\,)$\; Ricci\;
flow\; on\; a\; three--manifold\; which\; is\; complete\; with
bounded curvature for each $t\geq0$. Assume at $t=0$ the
eigenvalues $\lambda\geq\mu\geq\nu$ of the curvature operator at
each point are bounded below by $\nu\geq-1$. Then at all points
and all times $t\geq0$ we have the pinching estimate
$$
R\geq(-\nu)[\log(-\nu)+\log(1+t)-3]
$$
whenever $\nu<0$.
\end{theorem}
\begin{definition}[collapsed nonsingular normalized solution]
  \label{def:cz-collapsed-nonsingular-solution}
  \source{cao-zhu}
  \dcref{Definition 5.3.3}
  \group{Nonsingular Three-Dimensional Solutions}

We say a solution to the normalized Ricci flow is {\bf
collapsed}\index{collapsed} if there is a sequence of times
$t_k\rightarrow+\infty$ such that $\hat{\rho}(t_k)\rightarrow0$ as
$k\rightarrow+\infty$.
\end{definition}
\begin{theorem}[structure of noncollapsing nonsingular three-dimensional flows]
  \label{thm:cz-noncollapsing-nonsingular-flow-structure}
  \source{cao-zhu}
  \dcref{Theorem 5.3.4}
  \group{Nonsingular Three-Dimensional Solutions}

Let $g_{ij}(t)$, $0\leq t<+\infty$, be a noncollapsing nonsingular
solution of the normalized Ricci flow on a compact three-manifold
$M$. Then either
\begin{itemize}
\item[(i)] there exist a sequence of times $t_k\rightarrow+\infty$
and a sequence of diffeomorphisms $\varphi_k:\ M\rightarrow M$ so
that the pull-back of the metric $g_{ij}(t_k)$ by $\varphi_k$
converges in the $C^\infty$ topology to a metric on $M$ with
constant sectional curvature; or \item[(ii)] we can find a finite
collection of complete noncompact hyperbolic three-manifolds
$\mathcal{H}_1,\ldots,\mathcal{H}_m$ with finite volume, and for
all $t$ beyond some time $T<+\infty$ we can find compact subsets
$K_1,\ldots,K_m$ of $\mathcal{H}_1,\ldots,\mathcal{H}_m$
respectively obtained by truncating each cusp of the hyperbolic
manifolds along constant mean curvature torus of small area, and
diffeomorphisms $\varphi_l(t)$, $1\leq l\leq m$, of $K_l$ into $M$
so that as long as $t$ sufficiently large, the pull-back of the
solution metric $g_{ij}(t)$ by $\varphi_l(t)$ is as close as to
the hyperbolic metric as we like on the compact sets
$K_1,\ldots,K_m$; and moreover if we call the \textbf{exceptional
part} of $M$ those points where they are not in the image of any
$\varphi_l$, we can take the injectivity radii of the exceptional
part at everywhere as small as we like and the boundary tori of
each $K_l$ are \textbf{incompressible} in the sense that each
$\varphi_l$ injects $\pi_1(\partial K_l)$ into $\pi_1(M)$.
\index{exceptional part} \index{incompressible}
\end{itemize}
\end{theorem}
\begin{remark}[geometry of the exceptional long-time region]
  \label{rem:cz-exceptional-long-time-region-geometry}
  \source{cao-zhu}
  \dcref{Remark 5.3.5}
  \group{Nonsingular Three-Dimensional Solutions}

The exceptional part has bounded curvature and arbitrarily small
injectivity radii everywhere as $t$ large enough. Moreover the
boundary of the exceptional part consists of a finite disjoint
union of tori with sufficiently small area and is convex. Then by
the work of Cheeger-Gromov \cite{ChG86}, \cite{ChG90} or
Cheeger-Gromov-Fukaya \cite{CGF92}, there exists an
$\mathcal{F}$-structure on the exceptional part. In particular,
the exceptional part is a graph manifold, which have been
topologically classified. Hence \textbf{any nonsingular solution
to the normalized Ricci flow is geometrizable} in the sense of
Thurston (see the last section of Chapter 7 for details).
\end{remark}
\begin{lemma} [hyperbolicity of noncollapsing limits]
  \label{lem:cz-noncollapsing-limits-hyperbolic}
  \source{cao-zhu}
  \dcref{Lemma 5.3.6}
  \group{Nonsingular Three-Dimensional Solutions}

In Case $(3)$ where $R_{\min}\rightarrow-6$ as $t\rightarrow
+\infty$, all noncollapsing limit are hyperbolic with constant
sectional curvature $-1.$
\end{lemma}
\begin{lemma}[stability of finite-volume hyperbolic embeddings]
  \label{lem:cz-finite-volume-hyperbolic-embedding-stability}
  \source{cao-zhu}
  \dcref{Lemma 5.3.7}
  \group{Nonsingular Three-Dimensional Solutions}
 For any complete noncompact
hyperbolic three-manifold $\mathcal{H}$ with finite volume with
metric $h$, we can find a compact set $\mathcal{K}$ of
$\mathcal{H}$ such that for every integer $k$ and every
$\varepsilon>0$, there exist an integer $q$ and a $\delta>0$ with
the following property: if $F$ is a diffeomorphism of
$\mathcal{K}$ into another complete noncompact hyperbolic
three-manifold $\tilde{\mathcal{H}}$ with no fewer cusps $($than
$\mathcal{H}),$ finite volume with metric $\tilde{h}$ such that
$$
\|F^*\tilde{h}-h\|_{C^q(\mathcal{K})}<\delta
$$
then there exists an isometry $I$ of $\mathcal{H}$ to
$\tilde{\mathcal{H}}$ such that
$$
d_{C^k(\mathcal{K})}(F,I)<\varepsilon.
$$
\end{lemma}
\begin{lemma} [metric perturbations with concave boundary]
  \label{lem:cz-concave-boundary-metric-perturbations}
  \source{cao-zhu}
  \uses{lem:cz-hyperbolic-cusp-length-bounded}
  \dcref{Lemma 5.3.8}
  \group{Nonsingular Three-Dimensional Solutions}

Let\; $(X,\,g)$\; be\; a\; compact\; Riemannian\; manifold\; with
strictly negative Ricci curvature and with strictly concave
boundary. Then there are positive integer $l_0$ and small number
$\varepsilon_0>0$ such that for each positive integer $l\ge l_0$ and
positive number $\varepsilon\le\varepsilon_0$ we can find positive
integer $q$ and positive number $\delta>0$ such that for every
metric $\tilde{g}$ on $X$ with $||\tilde{g}-g||_{C^q(X)}\le \delta$
we can find a unique diffeomorphism $F$ of $X$ to itself so that
\begin{itemize}
\item[(a)] $F: (X,g)\rightarrow (X,\tilde{g})$ is harmonic,
\item[(b)] $F$ takes the boundary $\partial X$ to itself and
satisfies the {\bf free boundary condition}\index{free boundary
condition} that the normal derivative $\nabla_NF$ of $F$ at the
boundary is normal to the boundary, \item[(c)]
$d_{C^l(X)}(F,Id)<\varepsilon$, where $Id$ is the identity map.
\end{itemize}
\end{lemma}
\section{Named intermediate formalization obligations}
\begin{lemma}[curvature-gradient estimate under three-dimensional pinching]
  \label{lem:cz-three-dimensional-pinching-curvature-gradient}
  \source{cao-zhu}
  \dcref{Claim 1}
  \group{Differentiable Sphere Theorems}
For any $\varepsilon>0$, there exists a positive
constant $C_\varepsilon<+\infty$ such that for any time $\tau\ge
0$, we have
$$
\max_{t\le\tau}\max_{x\in M}|\nabla Rm(x,t)| \le \varepsilon
\max_{t\le\tau}\max_{x\in M}|Rm(x,t)|^\frac{3}{2}+C_\varepsilon.
$$

\end{lemma}
\begin{lemma}[bounded cusp length under nonnegative area derivative]
  \label{lem:cz-hyperbolic-cusp-length-bounded}
  \source{cao-zhu}
  \dcref{Claim 2}
  \group{Nonsingular Three-Dimensional Solutions}
If
$$
(1+2\varepsilon_0)L-(1-\varepsilon_0)A\geq 0
$$
then $L=L(\zeta)$ is uniformly bounded from above.

\end{lemma}
